\chapter*{Scope and assumptions}
\section*{General theory}
The general theory concerns finite transformation semigroups, their kernels,
partitions, ranges, idempotents, representations and associated operators.
It includes the partial-circulant family (PtC). Simplicity, the absence of
loops, two-out regularity, irreducibility and aperiodicity are not blanket
assumptions on this general theory. Probabilistic results state their
own hypotheses on the stochastic matrix and generator probabilities;
existence, positivity, uniqueness and convergence of stationary laws must
not be inferred without the relevant hypotheses.

\section*{Assumptions for road-coloring questions}
In questions concerning the road coloring problem, the underlying graph
$G$ is a finite simple directed graph: there are no loops and no parallel
edges with the same ordered endpoints. Opposite arrows $x\to y$ and
$y\to x$ are permitted. Every vertex has exactly two distinct outgoing
neighbours. The graph is strongly connected and aperiodic.
Write $\mathcal A$ for its adjacency matrix and set
\[
 \mathcal A\in\{0,1\}^{n\times n},\qquad
 \mathcal A_{xx}=0,\qquad\mathcal A\mathbf1=2\mathbf1,
 \qquad A=\tfrac12\mathcal A=\tfrac12(R+B).
\]
Thus $A$ is stochastic, traceless, irreducible and aperiodic. For a coloring,
\[
 R_{xx}=B_{xx}=0,\qquad R_{xy}B_{xy}=0,
 \qquad xR\ne xB.
\]
The stationary row $\pi$ is strictly positive and normalized by
$\pi\mathbf1=1$; $\pi A=\pi$ and $A^m\to\mathbf1\pi$.
Within this setting $A$ denotes the normalized transition matrix, not the
adjacency matrix. Recoloring changes $R,B$ but not $\mathcal A$ or $A$.

The auxiliary target-pair graph has one edge $\{xR,xB\}$ for each source.
Under these road-coloring assumptions it has no loops. Distinct sources
may give the same auxiliary edge without violating simplicity of $G$.

The road-coloring assumptions are not imposed on PtC or on the general
semigroup theory. Each theorem retains its own explicitly stated hypotheses.
